\section{CR manifolds and the Fefferman metric}\label{CR}
\subsection{CR structures}
An {\it almost CR structure} on a $(2n+1)$-dimensional $C^\infty$ manifold $M$ is defined to be a rank $n$ complex vector sub-bundle $T^{1, 0}M\subset\mathbb{C}TM$ satisfying $T^{1, 0}M\cap T^{0, 1}M=0$, where \smash{${T^{0, 1}M:=\ol{T^{1, 0}M}}$}. If the space of sections of $T^{1, 0}M$ is closed under the Lie bracket, it is called a {\it CR structure}. Since $T^{1, 0}M$ contains no non-zero real vector, the real part $H:=\operatorname{Re} T^{1, 0}M\subset TM$ becomes a~rank $2n$ real sub-bundle satisfying $\mathbb{C}H=T^{1, 0}M\oplus T^{0,1}M$.

We say a CR structure $T^{1, 0}M$ is {\it non-degenerate} if $H$ is a contact distribution, or equivalently, the {\it Levi form}
\[
T^{1, 0}_xM\times T^{0, 1}_xM\ni\bigl(Z, \ol W\bigr)\longmapsto {\rm i}\bigl[Z, \ol W\bigr]
\in T_xM/H_x
\]
gives a non-degenerate $TM/H$-valued hermitian form on \smash{$T_x^{1, 0}M$} at any point $x\in M$. We note that the right-hand side is well-defined via arbitrary extensions of $Z$, $\ol W$ to vector fields.

We usually assume that the real line bundle $TM/H$ is orientable so that we can trivialize $H^\perp \subset T^*M$ by a global contact form \smash{$\th\in \Gamma\bigl(H^\perp\bigr)$}. For a fixed orientation of $TM/H$, a contact form is determined up to multiplications by positive functions:
\smash{$\th\mapsto \wh\th={\rm e}^{\Upsilon}\th$, $\Upsilon\in C^\infty(M)$}. For each choice of $\th$, we have the real valued Levi form $h_\th$ by trivializing $TM/H$, which transforms conformally as $h_{\wh\th}={\rm e}^\Upsilon h_\th$.

The {\it Reeb vector field} of a contact form $\th$ is a real vector field $T$ uniquely
specified by the conditions $\th(T)=1$, $T\lrcorner\, {\rm d}\th=0$. Choosing a local frame $(Z_\a)$ for $T^{1, 0}M$, we have a local frame \smash{$\bigl(T, Z_\a, Z_{\ol\a}:=\ol{Z_\a}\bigr)$} for $\mathbb{C}TM=\mathbb{C}T\oplus T^{1, 0}M\oplus T^{0,1}M$ called an {\it admissible frame}, and its dual coframe \smash{$\bigl(\th, \th^\a, \th^{\ol\a}=\ol{\th^\a}\bigr)$}. In this coframe, we can write as
\[
{\rm d}\th={\rm i} h_{\a\ol\b}\th^\a\wedge\th^{\ol\b}
\]
and the Levi form is given by $\bigl(Z, \ol W\bigr)\mapsto h_{\a\ol\b}Z^\a \ol{W^\b}$.

%%%%%%%%%%%%%%%%%%%%%%%%%%%%%%%%%%%%%%%%%%%%%%%%%%%%%%%%%%%%%%%
\subsection{The Fefferman metric}\label{CR-Fefferman-metric}
The Fefferman metric associated to CR manifolds was first introduced by Fefferman~\cite{F} for the boundary of strictly pseudoconvex domains in $\mathbb{C}^{n+1}$ via an indefinite K\"ahler metric on $\mathbb{C}^{n+1}\times\mathbb{C}^*$ called the ambient metric. Burns--Diederich--Shnider~\cite{BDS} gave a construction via the CR Cartan bundle, and intrinsic characterizations were obtained by Farris~\cite{Fa} and Lee~\cite{Lee}.
Here, we~introduce the Fefferman metric by following Lee's formulation.

Let $\bigl(M, T^{1, 0}M\bigr)$ be a non-degenerate $(2n+1)$-dimensional CR manifold.
The {\it CR canonical bundle} is the complex line bundle over $M$ defined by
\[
K_M:=\biggl\{\xi\in\bigwedge\nolimits^{\!n+1}(\mathbb{C}T^*M)\, \Big|\, \ol Z\lrcorner\, \xi=0\ \textrm{for any}\ \ol Z\in T^{0,1}M\biggr\}.
\]
We note that a choice of an admissible frame gives a local frame $\th\wedge\th^1\wedge\cdots\wedge\th^n$ for $K_M$. The {\it Fefferman space} is the associated $S^1$-bundle
\[
\mathcal{C}:=K^\times_M/\mathbb{R}_+
\]
over $M$. If we choose a contact form $\th$, the Fefferman space is realized as a sub-bundle of $K^\times_M$~as%
\begin{equation}\label{embed-C}
\mathcal{C}\cong \bigl\{{\rm e}^{{\rm i}s}\th\wedge\th^1\wedge\cdots\wedge\th^n\mid s\in \mathbb{R}/2\pi\mathbb{Z}\bigr\} \subset K^\times_M,
\end{equation}
where $\bigl(\th, \th^\a, \th^{\ol\a}\bigr)$ is an admissible coframe such that $\bigl|\det\bigl(h_{\a\ol\b}\bigr)\bigr|=1$. This realization defines a~tautological $(n+1)$-form
\[
\xi={\rm e}^{{\rm i}s}\th\wedge\th^1\wedge\cdots\wedge\th^n
\]
on $\mathcal{C}$. Also, $\varphi={\rm e}^{{\rm i}s}\th^1\wedge\cdots\wedge\th^n$ becomes a global $n$-form on $\mathcal{C}$ uniquely characterized by the conditions
\[
\xi=\th\wedge\varphi, \qquad V\lrcorner\, \varphi=0\quad \text{for any lift}\ V\ \text{of}\ T;
\]
see {\cite[Lemma 3.3]{Lee}}.

We also have a $1$-form $\sigma$ on $\mathcal{C}$ specified by the following proposition.

\begin{Proposition}[{\cite[Proposition 3.4]{Lee}}]\label{sigma-condition}
There exists a unique real $1$-form $\sigma$ on $\mathcal{C}$ satisfying
\begin{align*}
&{\rm d}\xi={\rm i}(n+2)\sigma\wedge\xi, \\
&\sigma\wedge {\rm d}\varphi\wedge\ol\varphi=(\operatorname{Tr} {\rm d}\sigma){\rm i}\sigma\wedge\th\wedge\varphi\wedge\ol\varphi.
\end{align*}
\end{Proposition}
Here, for a 2-form $\omega$ on $M$ written as \smash{$\omega={\rm i}\omega_{\a\ol\b}\th^\a\wedge\th^{\ol\b}+\cdots$} in an admissible coframe, we define
\smash{$\operatorname{Tr} \omega:=h^{\a\ol\b}\omega_{\a\ol\b}$} using the inverse \smash{$h^{\a\ol\b}$} of the Levi form. The $1$-form $(n+2)\sigma$ gives a~principal connection form on the $S^1$-bundle $\mathcal{C}$.

By using $\sigma$, we now introduce the {\it Fefferman metric} as the pseudo-Riemannian metric on $\mathcal{C}$ defined by
\[
g^{\mathrm{F}}:=2h_{\a\ol\b}\theta^\a\cdot \theta^{\ol\b}+4\theta\cdot \sigma.
\]
We note that our definition of $g^{\mathrm{F}}$ is {\it twice} the definition of Lee. In the coframe $\smash{\bigl(\th^0:=\th, \th^\a, \th^{\ol\a}},\allowbreak \smash{\th^{n+1}:=2\sigma\bigr)}$, the non-zero components of the metric tensor are
\[
g^{\mathrm{F}}_{\a\ol\b}=g^{\mathrm{F}}_{\ol\b\a}=h_{\a\ol\b}, \qquad
g^{\mathrm{F}}_{0\, n+1}=g^{\mathrm{F}}_{n+1\, 0}=1.
\]
If the Levi form has signature $(p, q)$ as a hermitian form, then the Fefferman metric has signature $(2p+1, 2q+1)$. Moreover, if we rescale the contact form
as $\wh\th={\rm e}^\Upsilon\th$, the Fefferman metric transforms as
$\wh{g}^{\,\mathrm{F}}={\rm e}^\Upsilon g^{\mathrm{F}}$. Thus, the conformal structure $\bigl[g^{\mathrm{F}}\bigr]$ is independent of the choice of contact form~\cite[Theorem~5.17]{Lee}.

It follows from the formula in~\cite{Lee} of the connection form of $g^{\mathrm{F}}$ in terms of the pseudo-hermitian structure of $M$ that $g^{\mathrm{F}}$ is conformally flat if and only if $M$ is locally CR equivalent to the hyperquadric
\[
\mathcal{Q}=\bigl\{[\xi]\in \mathbb{CP}^{n+1} \,\big|\, -\bigl|\xi^0\bigr|^2-\bigl|\xi^1\bigr|^2-\cdots-
\bigl|\xi^q\bigr|^2+\bigl|\xi^{q+1}\bigr|^2+\cdots+\bigl|\xi^{n+1}\bigr|^2=0 \bigr\},
\]
which is the flat model of CR manifold.

\subsection{Chains and null chains}
It is well-known that null geodesics of a pseudo-Riemannian metric are conformally invariant up to reparametrizations. Hence, the projections of null geodesics of the Fefferman metric provide an invariant family of curves on CR manifolds. Let $\g(t)$ be a null geodesic of $g^{\mathrm{F}}$. Since $g^{\mathrm{F}}$ is $S^1$-invariant, the infinitesimal generator $K$ of the $S^1$-action on $\mathcal{C}$ is a Killing vector field and hence $g^{\mathrm{F}}(K, \dot\g(t))$ is constant in $t$.

\begin{Definition}[\cite{F, Koch}]
An unparametrized curve on $M$ is called a {\it chain} (resp.\ {\it null chain}) if it is the projection of a non-vertical null geodesic $\g(t)$ of $\bigl[g^{\mathrm{F}}\bigr]$ with $g^{\mathrm{F}}(K, \dot\g)\neq0$ \big(resp.\ $g^{\mathrm{F}}(K, \dot\g)=0$\big).
\end{Definition}

Since $K\lrcorner\, g^{\mathrm{F}}=2\th$, the definition implies that chains are transverse to the contact distribution~$H$ while null chains are tangent to $H$ and null with respect to the Levi form.

For each choice of a direction transverse to $H_x$, there exists a unique chain emanating in that direction. Thus, chains can be considered as `geodesics' in CR geometry. In fact, they are characterized as geodesics of a Kropina metric on CR manifolds, which is a Finsler metric singular along $H$; see~\cite{CMMM} and Section~\ref{Kropina} below. On the other hand, there exists a 1-parameter family of null chains which are tangent to a fixed null vector in $H_x$; see~\cite{Koch}.

%%%%%%%%%%%%%%%%%%%%%%%%%%%%%%%%%%%%%%%%%%%%%%%%%%%%%%%%%%%%%%%%%
\section{Conformal three-manifolds and twistor CR manifolds}\label{conformal-three}

\subsection{Conformal three-manifolds}
Let $(\Sigma, [g])$ be a $3$-dimensional conformal manifold (of Riemannian signature). We take a representative metric $g\in[g]$ and let $R_{ij}{}^k{}_l$ be the curvature tensor of $g$, defined by $(\nabla_i\nabla_j-\nabla_j\nabla_i)X^k=R_{ij}{}^k{}_lX^l$. Since the Weyl curvature always vanishes in three dimensions, we can write as
\begin{equation}\label{curvature-tensor}
R_{ijkl}=P_{ik}g_{jl}-P_{jk}g_{il}+P_{jl}g_{ik}-P_{il}g_{jk}
\end{equation}
with the Schouten tensor
\[
P_{ij}:=R_{ij}-\frac{1}{4}Rg_{ij}.
\]
Here $R_{ij}:=R_{ki}{}^k{}_j$ is the Ricci tensor and $R:=R_k{}^k$ is the scalar curvature.

The Cotton tensor $C=(C_{ijk})\in \Gamma\bigl(T^*\Sigma\otimes \wedge^2 T^*\Sigma\bigr)$ is defined by
\[
C(X, Y, Z):=(\nabla_Z P)(X, Y)-(\nabla_Y P)(X, Z), \qquad \text{i.e.,}\qquad
C_{ijk}=\nabla_k P_{ij}-\nabla_j P_{ik}.
\]
Since $\Sigma$ is 3-dimensional, the Cotton tensor is conformally invariant, and identically vanishes if and only if $(\Sigma, [g])$ is conformally flat.

If $\Sigma$ is oriented, we have the cross product $Z=X\times Y\in T_p\Sigma$ of tangent vectors $X, Y\in T_p\Sigma$. In the index notation, it is given by
\[
Z^i=\e^i{}_{jk}X^j Y^k,
\]
where $\e_{ijk}=\e_{[ijk]}$ is the volume form of $g$:
\[
vol_g=\frac{1}{3!}\e_{ijk}\th^i\wedge \th^j\wedge \th^k.
\]
Applying the Hodge star operator to the Cotton tensor, we define $*C\in\Gamma(T^*\Sigma\otimes T^*\Sigma)$ by
\[
(*C)_{ij}:=\frac{1}{2}\varepsilon_j{}^{kl}C_{ikl},
\]
which is characterized by the equation
\[
C(X, Y, Z)=(*C)(X, Y\times Z).
\]
Since $C_{ijk}$ is trace-free by the contracted second Bianchi identity, $(*C)_{ij}$ becomes a symmetric 2-tensor.

\subsection{Conformal geodesics}
Although geodesics of a Riemannian metric are not conformally invariant, there is an alternative conformally invariant family of curves on conformal manifolds, called {\it conformal geodesics}. They are characterized by the conformally invariant third order ODE
\begin{equation}\label{conf-geod-proj}
\nabla_{\dot x}\nabla_{\dot x}\dot x-\frac{3g(\dot x, \nabla_{\dot x}\dot x)}{|\dot x|^2}\nabla_{\dot x}\dot x+\frac{3|\nabla_{\dot x}\dot x|^2}{2|\dot x|^2}\dot x
+2P(\dot x, \dot x)\dot x-|\dot x|^2P(\dot x)=0.
\end{equation}
Here, we set $P(\dot x, \dot x):=P_{ij}\dot x^i\dot x^j$ and \smash{$P(\dot x)^i:=P^i{}_j \dot x^j$}. We refer the reader to~\cite{BEG} for a derivation of this equation by using the tractor calculus, which is an invariant calculus in conformal geometry.
Conformal geodesics are also called {\it conformal circles} since these curves are circles and straight lines when $g$ is the Euclidean metric.

It is shown in~\cite{BE} that if $x(t)$ is a solution to the equation \eqref{conf-geod-proj}, a reparametrized curve $\tilde x(t)=x(\varphi(t))$ solves \eqref{conf-geod-proj} if and only if $\varphi(t)$ is a projective linear transformation: $\varphi(t)=(at+b)/(ct+d)$. Hence, we call a solution to \eqref{conf-geod-proj}
a {\it conformal geodesic with a projective parameter}. In~\cite{BE}, they also prove that one can define an equation for `unparametrized conformal geodesics' by taking the $\dot x$-orthogonal part of \eqref{conf-geod-proj}.
\begin{Proposition}[{\cite[Proposition 4.2]{BE}}]
A regular curve $x(t)$ can be reparametrized to satisfy~\eqref{conf-geod-proj} if and only if it satisfies the equation
\begin{equation}\label{conf-geod}
\dot x\wedge\biggl(\nabla_{\dot x}\nabla_{\dot x}\dot x-\frac{3g(\dot x, \nabla_{\dot x}\dot x)}{|\dot x|^2}\nabla_{\dot x}\dot x-|\dot x|^2P(\dot x)\biggr)=0.
\end{equation}
\end{Proposition}
In this paper, we always consider conformal geodesics as unparametrized curves, and we say a~regular curve $x(t)$ is a {\it conformal geodesic} if it satisfies \eqref{conf-geod}.
We note that if $x(t)$ has a constant speed $|\dot x|$, the equation \eqref{conf-geod} is reduced to
\begin{equation}\label{conf-geod-const-speed}
\dot x\wedge\bigl(\nabla_{\dot x}\nabla_{\dot x}\dot x-|\dot x|^2P(\dot x)\bigr)=0
\end{equation}
since $g(\dot x, \nabla_{\dot x}\dot x)=0$.

Conformal geodesics are also characterized as follows.

\begin{Proposition}[{\cite[Proposition 3.3]{BE}}]\label{conf-geod-characterization}
A regular curve $x(t)$ is a conformal geodesic if and only if there exist a metric $g\in[g]$ and a reparametrization such that
\[
\nabla_{\dot x}\dot x=P(\dot x)=0.
\]
\end{Proposition}

\subsection{Twistor CR manifolds}
We will define the twistor CR manifold $M$ over $\bigl(\Sigma^3, [g]\bigr)$, which was introduced by LeBrun~\cite{LeB} as a $3$-dimensional analogue of Penrose's twistor space for anti-self dual conformal $4$-manifolds.

We extend a representative metric $g\in[g]$ on $T^*\Sigma$ to a complex bi-linear form on $\mathbb{C}T^*\Sigma$ and define a $7$-dimensional manifold $\hat M$ as the set of complex null covectors:
\[
\hat M:=\{\zeta\in\mathbb{C}T^*\Sigma\mid g(\zeta, \zeta)=0,\, \zeta\neq0\}.
\]
By taking the quotient by the $\mathbb{C}^*$-action, we obtain a $5$-dimensional manifold
\[
M:=\hat M/\mathbb{C}^*\subset \mathbb{P}(\mathbb{C}T^*\Sigma).
\]
We denote the natural projections by $\tilde\pi\colon \hat M\to M$, $\pi\colon M\to \Sigma$. Note that the fiber of $\pi$ is biholomorphic to \smash{$\bigl\{[\zeta]\in\mathbb{CP}^2\mid \zeta_1^2+\zeta_2^2+\zeta_3^2=0\bigr\}\cong \mathbb{CP}^1$}.

Let
\smash{$\alpha:=\bigl(\zeta_i {\rm d}x^i\bigr)|_{T\hat M}$} be the canonical (tautological) $1$-form on
$\hat M$ and set $\omega:={\rm d}\alpha$. Then, the distribution
\[
 D:=\textrm{Ker}\, \omega=\bigl\{v\in \mathbb{C}T\hat M\mid v\lrcorner\, \omega=0\bigr\}
\]
is $\mathbb{C}^*$-invariant and we can define the distribution $\tilde\pi_* D\subset \mathbb{C}TM$. LeBrun~\cite{LeB} proved that
\[
T^{1, 0}M:=\ol{\tilde\pi_* D}
\]
becomes a non-degenerate CR structure on $M$ whose Levi form has signature $(1, 1)$. The construction depends only on the conformal class $[g]$, and $\bigl(M, T^{1, 0}M\bigr)$ is called the {\it twistor CR manifold} over $(\Sigma, [g])$.

The CR structure $T^{1, 0}M$ can be described in terms of coordinates as follows.
Fix a point $x_0\in \Sigma$ and let $(x^i)$ be normal coordinates centered at $x_0$ with respect to a representative metric $g\in[g]$.
\begin{Proposition}\label{D}
For any point $[\zeta]\in \pi^{-1}(x_0)$, we have
\[
\ol D_{\zeta}=\mathbb{C}\ol{\zeta^i}\frac{\partial}{\partial x^i}\oplus \biggl\{ w_i\frac{\partial}{\partial \zeta_i} \, \bigg|\, \zeta_i w^i=0\biggr\},
\]
and the projection \smash{$\tilde\pi_*\colon \ol D_{\zeta}\to T^{1. 0}_{[\zeta]} M$} induces the isomorphism
\[
T^{1, 0}_{[\zeta]} M\cong\mathbb{C}\ol{\zeta^i}\frac{\partial}{\partial x^i}\oplus \biggl(\biggl\{ w_i\frac{\partial}{\partial \zeta_i}\, \bigg|\, \zeta_i w^i=0\biggr\}\Big/\mathbb{C}\zeta_i\frac{\partial}{\partial \zeta_i}\biggr).
\]
\end{Proposition}
See \cite[Proposition 6.1]{Mar} for the proof.

\section[The Fefferman metric and (null) chains of twistor CR manifolds]{The Fefferman metric and (null) chains\\ of twistor CR manifolds}\label{Fefferman-twistor-CR}
Let $\pi\colon M\to \Sigma$ be the twistor CR manifold over a conformal $3$-manifold $(\Sigma, [g])$. We will describe the Fefferman metric for $M$ and compute its geometric quantities.

We fix a representative metric $g\in[g]$ and define
\[
\wh{\mathcal{C}}:=\bigl\{\zeta\in \hat M\mid g\bigl(\zeta, \ol\zeta\bigr)=1\bigr\}
=\bigl\{\zeta\in\mathbb{C}T^*\Sigma\mid g(\zeta, \zeta)=0,\, g\bigl(\zeta, \ol\zeta\bigr)=1\bigr\},
\]
which is an $S^1$-bundle over $M$: \smash{$\wh{\mathcal{C}}/\mathrm{U}(1)=M$}.
We show that \smash{$\wh{\mathcal{C}}$} gives a double covering of the Fefferman space $\mathcal{C}$ and construct the Fefferman metric on \smash{$\wh{\mathcal{C}}$} instead of $\mathcal{C}$.



\subsection[An adapted frame on wh{C}]{An adapted frame on $\boldsymbol{\wh{\mathcal{C}}}$}
Let $(\zeta_i)$ be the fiber coordinates for $\zeta_i{\rm d}x^i\in \mathbb{C}T^*\Sigma$ associated to local coordinates $(x^i)$ on $\Sigma$. On~$\wh{\mathcal{C}}$, it holds that
\[
\zeta^i\zeta_i=0, \qquad \zeta^i\ol\zeta_i=1.
\]
By using the metric, we identify $\zeta_i {\rm d}x^i$ with the complex tangent vector
\[
\zeta=\zeta^i\frac{\partial}{\partial x^i}
\]
and regard that $\zeta$ also represents the global section of the vector bundle
$\hat\pi^{-1}(\mathbb{C}T\Sigma)$ over $\wh{\mathcal{C}}$, where \smash{$\hat\pi\colon \wh{\mathcal{C}}\to \Sigma$} is the projection. We (locally) choose an orientation of $\Sigma$ and define a real vector
\[
\eta=\eta^i\frac{\partial}{\partial x^i}\in \hat\pi^{-1}(T\Sigma)
\]
by
\[
\eta:={\rm i}\zeta\times \ol\zeta=2\operatorname{Re}\zeta\times \operatorname{Im}\zeta.
\]
Using the volume form $\e_{ijk}$ of $g$, we can write as
\[
\eta^j={\rm i}\e^{jkl}\zeta_k\ol{\zeta_l}.
\]
The property of the cross product gives
\begin{alignat*}{3}
&g(\eta, \eta)=1, \qquad&& g(\eta, \zeta)=g\bigl(\eta, \ol\zeta\bigr)=0,& \\
&\eta\times\ol\zeta={\rm i}\ol\zeta, \qquad&& \zeta\times\eta={\rm i}\zeta.&
\end{alignat*}
(In the index notation, one can use the identity $\e_{ijk}\e^i{}_{lm}=g_{jl}g_{km}-g_{kl}g_{jm}$ to verify these equations.)

Since the conditions $g(\zeta, \zeta)=0$, $g\bigl(\zeta, \ol\zeta\bigr)=1$ are preserved by the parallel transport of $\zeta$, we can define the horizontal lifts of tangent vectors on $\Sigma$ to \smash{$\wh{\mathcal{C}}$}. Let \smash{$\Gamma_{ij}{}^k$} be the Christoffel symbols of the Levi-Civita connection of $g$. Then, the horizontal lift of the tangent vector $\partial/\partial x^i$ is given~by
\[
\frac{\d}{\d x^i}:=\frac{\partial}{\partial x^i}+\Gamma_{ij}{}^k\zeta_k\frac{\partial}{\partial \zeta_j}+\Gamma_{ij}{}^k\ol{\zeta_k}\frac{\partial}{\partial \ol{\zeta}_j}\in T\wh{\mathcal{C}}.
\]
Note that this vector transforms tensorially in the index $i$ under coordinate changes. The annihilators of the horizontal distribution on $\wh{\mathcal{C}}$ are given by
\[
\d \zeta_i:={\rm d}\zeta_i-\Gamma_{ki}{}^j \zeta_j {\rm d}x^k, \qquad \d \ol{\zeta}_i:={\rm d}\ol{\zeta}_i-\Gamma_{ki}{}^j \ol{\zeta}_j {\rm d}x^k,
\]
which also transform tensorially.

Using these notations, we define a frame
\smash{$\bigl(\tilde T, \tilde Z_1, \tilde Z_2, \tilde Z_{\ol1}, \tilde Z_{\ol2}, K\bigr)$}
 for $\mathbb{C}T\wh{\mathcal{C}}$ and its dual coframe \smash{$\bigl(\th, \tilde\th^1, \tilde\th^2, \tilde\th^{\ol1}, \tilde\th^{\ol2}, \tilde\th^3\bigr)$} by
\[
\tilde T:=\eta^i \frac{\d}{\d x^i}, \qquad \tilde Z_1:=\ol{\zeta^i}\frac{\d}{\d x^i}, \qquad \tilde Z_2:=-{\rm i}\eta_j\frac{\partial}{\partial \zeta_j}, \qquad
K:={\rm i}\biggl(\zeta_j\frac{\partial}{\partial \zeta_j}
-\ol{\zeta}_j\frac{\partial}{\partial \ol{\zeta}_j}\biggr)
\]
and
\[
\th:=\eta_i {\rm d}x^i, \qquad \tilde\th^1:=\zeta_i {\rm d}x^i, \qquad \tilde\th^2:={\rm i}\eta^j\d\zeta_j, \qquad \tilde\th^3:=-\frac{{\rm i}}{2}\bigl(\ol{\zeta^j}\d\zeta_j-\zeta^j\d\ol{\zeta}_j\bigr)=-i\ol{\zeta^j}\d\zeta_j.
\]
Note that these vectors and forms are independent of the choice of coordinates and globally defined if $\Sigma$ is globally oriented. We also note that $K$ agrees with the infinitesimal generator of the $S^1$-action on $\wh{\mathcal{C}}$.

We shall prove that this indeed gives a (co)frame which is `adapted' to the CR structure on~$M$.

\begin{Proposition}\quad
\begin{itemize}\itemsep=0pt
\item[$(1)$] The $1$-from $\th$ descends to a contact form on $M$.
\item[$(2)$] The $1$-forms $\tilde\th^1$, $\tilde\th^2$ are horizontal with respect to $\wh{\mathcal{C}}\to M$, and for any local section ${s\colon M\to \wh{\mathcal{C}}}$, the set of $1$-forms $\bigl(\th, \th^1=s^*\tilde\th^1, \th^2=s^*\tilde\th^2, \th^{\ol1}=\ol{\th^1}, \th^{\ol2}=\ol{\th^2}\bigr)$ gives an admissible coframe on $M$.
\end{itemize}
\end{Proposition}
\begin{proof}
We take normal coordinates $(x^i)$ around a fixed point $x_0\in \Sigma$ and work at a point $p\in\pi^{-1}(x_0)\subset M$.

(1) Since each coefficient $\eta_i$ is a $\mathrm{U}(1)$-invariant function, $\th$ descends to a $1$-form on $M$. Moreover, by Proposition~\ref{D}, $\th$ annihilates $T^{1, 0}M\oplus T^{0, 1}M$ hence is a contact form on $M$. We note that the projection $T$ of $\tilde T$ to $M$ is well-defined since $\tilde T$ is $\mathrm{U}(1)$-invariant, and it gives the Reeb vector field of $\th$. In fact, $T$ is represented as $T=\eta^i\bigl(\partial/\partial x^i\bigr)$ at $p$, and since
${\rm d}\e_{ijk}(x_0)=0$, ${\rm d}g_{ij}(x_0)=0$, we have ${\rm d}\eta_i(T)=0$, $\eta^i{\rm d}\eta_i=0$ at $p$ and hence $\th(T)=\eta_i\eta^i=1$, $T\lrcorner\, {\rm d}\th=T\lrcorner\, \bigl({\rm d}\eta_i\wedge {\rm d}x^i\bigr)=0$.

(2) Since $\tilde\th^1(K)=\tilde\th^2(K)=0$, these forms are horizontal. By Proposition~\ref{D}, we see that~$\th^1$ and~$\th^2$ annihilate $T^{0, 1}M$. Moreover, $\tilde\th^1\bigl(\tilde T\bigr)=\tilde\th^2\bigl(\tilde T\bigr)=0$ implies \smash{$\th^1(T)=\th^2(T)=0$}. Thus, \smash{$\bigl(\th, \th^1, \th^2, \th^{\ol1}, \th^{\ol2}\bigr)$} is an admissible coframe.
\end{proof}

We remark that $\tilde\th^1$ and $\tilde\th^2$ do not descend to $1$-forms on $M$. In fact, if we denote the action of $\lambda\in \mathrm{U}(1)$ on $\wh{\mathcal{C}}$ by $\d_\lambda\colon \wh{\mathcal{C}}\to \wh{\mathcal{C}}$, we have
\[
\d_\lambda^* \tilde\th^1=\lambda \tilde\th^1, \qquad \d_\lambda^*\tilde\th^2=\lambda\tilde\th^2.
\]
Thus, $\bigl(\tilde\th^1, \tilde\th^2\bigr)$ can be considered as a `twisted' or `weighted' coframe for $T^{1, 0}M$.

We also note that if we replace $g$ with $\wh g={\rm e}^{2\Upsilon} g$, $\Upsilon\in C^\infty(\Sigma)$, we have $\wh\eta^{\,i}={\rm e}^{-\Upsilon}\eta^i$, $\wh\eta_i={\rm e}^{\Upsilon}\eta_i$ and hence the associated contact form is rescaled as $\wh\th={\rm e}^{\Upsilon}\th$.

%%%%%%%%%%%%%%%%%%%%%%%%%%%%%%%%%%%%%%%%%%%%%%%%%%%%%%

\subsection{The differentials of the adapted coframe}
We will compute the exterior derivatives of the adapted coframe defined above.
We use the following lemma in computations in normal coordinates.
\begin{Lemma}
Let $(x^i)$ be normal coordinates of $g$ centered at a point $x_0\in\Sigma$. Then, we have
\begin{equation}\label{d-eta}
{\rm d}x^j=\eta^j\theta+\ol{\zeta^j}\tilde\theta^1+\zeta^j\tilde\theta^{\ol1}, \qquad
{\rm d}\zeta_j=-{\rm i}\bigl(\eta_j\tilde\theta^2-\zeta_j\tilde\theta^3\bigr), \qquad
{\rm d}\eta^j={\rm i}\bigl(\ol{\zeta^j}\tilde\theta^2-\zeta^j\tilde\theta^{\ol2}\bigr)
\end{equation}
at any point $p\in \pi^{-1}(x_0)\subset M$.
\end{Lemma}
\begin{proof}
At $p$, we have
\[
{\rm d}x^j\bigl(\tilde T\bigr)=\eta^j, \qquad {\rm d}x^j\bigl(\tilde Z_1\bigr)=\ol{\zeta^j},\qquad {\rm d}x^j\bigl(\tilde Z_2\bigr)={\rm d}x^j(K)=0,
\]
from which we obtain the first equation in \eqref{d-eta}. Similarly, we have
\[
{\rm d}\zeta_j\bigl(\tilde T\bigr)={\rm d}\zeta_j\bigl(\tilde Z_{1}\bigr)={\rm d}\zeta_j\bigl(\tilde Z_{\ol1}\bigr)={\rm d}\zeta_j\bigl(\tilde Z_{\ol2}\bigr)=0, \qquad
{\rm d}\zeta_j\bigl(\tilde Z_2\bigr)=-{\rm i}\eta_j, \qquad
{\rm d}\zeta_j(K)={\rm i}\zeta_j,
\]
and obtain the formula for ${\rm d}\zeta_j$. Finally, using $d\e_{ijk}(x_0)=0$, we compute as
 \[
{\rm d}\eta={\rm i}\,{\rm d}\zeta\times \ol\zeta+{\rm i}\zeta\times {\rm d}\ol\zeta
=\bigl(\eta\tilde\theta^2-\zeta\tilde\theta^3\bigr)\times\ol\zeta-\zeta\times\bigl(\eta\tilde\theta^{\ol2}-\ol\zeta\tilde\theta^3\bigr)={\rm i}\ol\zeta\tilde\theta^2-{\rm i}\zeta\tilde\theta^{\ol2}
\]
at $p$.
\end{proof}

\begin{Proposition} We have
\begin{gather}
{\rm d}\theta={\rm i}\bigl(\tilde\theta^1\wedge\tilde\theta^{\ol2}+\tilde\theta^2\wedge\tilde\theta^{\ol1}\bigr), \nonumber\\
{\rm d}\tilde\theta^1={\rm i}\theta\wedge\tilde\theta^2-{\rm i}\tilde\theta^1\wedge\tilde\theta^3, \nonumber\\
{\rm d}\tilde\theta^2={\rm i}\bigl(P\bigl(\zeta, \ol\zeta\bigr)+P(\eta, \eta)\bigr)\theta\wedge\tilde\theta^1+{\rm i}P(\zeta, \zeta)\theta\wedge\tilde\theta^{\ol1}
- {\rm i}P(\zeta, \eta)\tilde\theta^1\wedge\tilde\theta^{\ol1}-{\rm i}\tilde\theta^2\wedge\tilde\theta^3, \nonumber\\
{\rm d}\tilde\theta^3=- {\rm i}P\bigl(\ol\zeta, \eta\bigr)\theta\wedge\tilde\theta^1+{\rm i}
P(\zeta, \eta)\theta\wedge\tilde\theta^{\ol1}
+ 2{\rm i}P\bigl(\zeta, \ol\zeta\bigr)\tilde\theta^1\wedge\tilde\theta^{\ol1}+{\rm i}\tilde\theta^2\wedge\tilde\theta^{\ol2}.\label{diff-coframe}
\end{gather}
\end{Proposition}
\begin{proof}
We compute in normal coordinates around $x_0\in\Sigma$. By \eqref{d-eta}, we have
\begin{align*}
&{\rm d}\theta={\rm d}\eta_i\wedge {\rm d}x^i={\rm i}\bigl(\tilde\theta^1\wedge\tilde\theta^{\ol2}+\tilde\theta^2\wedge\tilde\theta^{\ol1}\bigr), \\
&{\rm d}\tilde\th^1={\rm d}\zeta_i\wedge {\rm d}x^i={\rm i}\theta\wedge\tilde\theta^2-{\rm i}\tilde\theta^1\wedge\tilde\theta^3.
\end{align*}
To compute ${\rm d}\tilde\theta^2$, we note that
\begin{align*}
{\rm d}\bigl(\Gamma_{kjl} {\rm d}x^k\bigr)&=\frac{1}{2}R_{kmlj}{\rm d}x^k\wedge {\rm d}x^m \\
&=\frac{1}{2}(P_{kl}g_{mj}-P_{ml}g_{kj}+P_{mj}g_{kl}-P_{kj}g_{ml}){\rm d}x^k\wedge {\rm d}x^m \\
&=(P_{kl}g_{mj}+P_{mj}g_{kl}){\rm d}x^k\wedge {\rm d}x^m
\end{align*}
holds at $p\in \pi^{-1}(x_0)$ since $\Gamma_{ij}{}^k(x_0)=0$. Using this equation and \eqref{d-eta}, we have
\begin{align*}
{\rm d}\tilde\theta^2={}&{\rm i}\,{\rm d}\eta^j\wedge {\rm d}\zeta_j-{\rm i}\eta^j\zeta^l{\rm d}\bigl(\Gamma_{kjl} {\rm d}x^k\bigr) \\
={}&{\rm i}\,{\rm d}\eta^j\wedge {\rm d}\zeta_j-{\rm i}\eta^j\zeta^l (P_{kl}g_{mj}+P_{mj}g_{kl}){\rm d}x^k\wedge {\rm d}x^m \\
={}&{\rm i}\bigl(\ol{\zeta^j}\tilde\th^2-\zeta^j\tilde\th^{\ol2}\bigr)\wedge\bigl(\eta_j\tilde\theta^2-\zeta_j\tilde\theta^3\bigr) \\
&{-}\,{\rm i}\eta^j\zeta^l (P_{kl}g_{mj}+P_{mj}g_{kl})\bigl(\eta^k\theta+\ol{\zeta^k}\tilde\theta^1+\zeta^k\tilde\theta^{\ol1}\bigr)\wedge \bigl(\eta^m\theta+\ol{\zeta^m}\tilde\theta^1+\zeta^m\tilde\theta^{\ol1}\bigr) \\
={}&{\rm i}\bigl(P\bigl(\zeta, \ol\zeta\bigr)+P(\eta, \eta)\bigr)\theta\wedge\tilde\theta^1+{\rm i}P(\zeta, \zeta)\theta\wedge\tilde\theta^{\ol1} -{\rm i}P(\zeta, \eta)\tilde\theta^1\wedge\tilde\theta^{\ol1}-{\rm i}\tilde\theta^2\wedge\tilde\theta^3.
\end{align*}
The computation of ${\rm d}\tilde\th^3$ is similar, and we omit it.
\end{proof}

It follows from the first equation in \eqref{diff-coframe} that
\[
{\rm d}\th={\rm i}\bigl(\theta^1\wedge\theta^{\ol2}+\theta^2\wedge\theta^{\ol1}\bigr)
\]
holds for the admissible coframe $\th^1=s^*\tilde \th^1$, $\th^2=s^*\tilde \th^2$ determined by a local section $s\colon M\to \wh{\mathcal{C}}$. Hence the Levi form is given by
\[
\bigl(h_{\a\ol\b}\bigr)=\begin{pmatrix} 0 & 1 \\ 1 & 0 \end{pmatrix},
\]
which indeed has signature $(1, 1)$.

\subsection{The Fefferman space for twistor CR manifolds}
Let $\mathcal{C}=K_M^\times/\mathbb{R}_+$ be the Fefferman space over $M$. For simplicity, we assume that $\Sigma$ is globally oriented, and let $\th=\eta_i{\rm d}x^i$ be the contact form determined by a choice of representative metric $g\in[g]$ as above.
We embed $\mathcal{C}$ into $K^\times_M$ as in \eqref{embed-C} by $\th$, and define the mapping
\[
\varpi\colon \ \wh{\mathcal{C}}\longrightarrow \mathcal{C}, \qquad u\longmapsto \bigl(\theta\wedge\tilde\theta^1\wedge\tilde\theta^2\bigr)_u
\]
by using our twisted coframe $\bigl(\tilde\th^1, \tilde\th^2\bigr)$ for $T^{1, 0}M$. %with $|\det(h_{\a\ol\b})|=1$.
Since we have
\[
\d_\lambda^*\bigl(\theta\wedge\tilde\theta^1\wedge\tilde\theta^2\bigr)
=\lambda^2\bigl(\theta\wedge\tilde\theta^1\wedge\tilde\theta^2\bigr)
\]
for the action of $\lambda\in\mathbb{C}^*$, the mapping $\varpi$ becomes a double covering map.

We will compute the pullback $\varpi^*g^{\mathrm{F}}$ of the Fefferman metric, which we also denote by $g^{\mathrm{F}}$ and call the {\it Fefferman metric} on $\wh{\mathcal{C}}$. We set
\[
\hat{\xi}:=\varpi^*\xi=\theta\wedge\tilde\theta^1\wedge\tilde\theta^2,\qquad \hat{\varphi}:=\varpi^*\varphi, \qquad \hat{\sigma}:=\varpi^*\sigma,
\]
where $\xi$, $\varphi$, $\sigma$ are the differential forms defined in Section~\ref{CR-Fefferman-metric}.
\begin{Lemma}
We have $\hat\varphi=\tilde\theta^1\wedge\tilde\theta^2$ and $\hat\sigma=\frac{1}{2}\tilde\th^3$.
\end{Lemma}
\begin{proof}
The 2-form $\hat\varphi=\tilde\theta^1\wedge\tilde\theta^2$ satisfies
$\hat\xi=\th\wedge\hat\varphi$ and $V\lrcorner\,\hat\varphi=0$ for any lift $V$ of $T$ since $\tilde T\lrcorner\,\hat\varphi= K\lrcorner\,\hat\varphi=0$. Thus, it satisfies the required characterization. By Proposition~\ref{sigma-condition}, $\hat\sigma$ is characterized by the conditions
\begin{align*}
&{\rm d}\hat\xi=4{\rm i}\hat\sigma\wedge\hat\xi, \\
&\hat\sigma\wedge {\rm d}\hat\varphi\wedge\ol{\hat\varphi}=(\operatorname{Tr} {\rm d}\hat\sigma){\rm i}\hat\sigma\wedge\th\wedge\hat\varphi\wedge\ol{\hat\varphi}.
\end{align*}
By \eqref{diff-coframe}, we have
\[
{\rm d}\hat\xi=2{\rm i}\tilde\theta^3\wedge\hat\xi, \qquad
\tilde\theta^3\wedge {\rm d}\hat\varphi\wedge \ol{\hat\varphi}=0.
\]
Moreover, since
\begin{align*}
{\rm d}\tilde\theta^3 =2{\rm i}P\bigl(\zeta, \ol\zeta\bigr)\tilde\theta^1\wedge\tilde\theta^{\ol1}+{\rm i}\tilde\theta^2\wedge\tilde\theta^{\ol2}+\cdots
 =2{\rm i}P\bigl(\zeta, \ol\zeta\bigr)\theta^1\wedge\theta^{\ol1}+{\rm i}\theta^2\wedge\theta^{\ol2}+\cdots,
\end{align*}
we have $\operatorname{Tr}{\rm d}\tilde\theta^3=0$. Thus, we obtain \smash{$\hat\sigma=\frac{1}{2}\tilde\theta^3$}.
\end{proof}

Since \smash{$g^{\mathrm{F}}=2h_{\a\ol\b}\theta^\a\cdot \theta^{\ol\b}+4\theta\cdot \hat\sigma=2h_{\a\ol\b}\tilde\theta^\a\cdot \tilde\theta^{\ol\b}+4\theta\cdot \hat\sigma$}, we obtain the following theorem.

\begin{Theorem}\label{Fefferman-metric}
The Fefferman metric on the double covering $\wh{\mathcal{C}}$ %determined by a representative metric $g\in[g]$
is given by
\[
g^{\mathrm{F}}=2\bigl(\tilde\theta^1\cdot\tilde\theta^{\ol2}+\tilde\theta^2\cdot\tilde\theta^{\ol1}+\theta\cdot \tilde\theta^3\bigr)
\]
in the adapted coframe.
\end{Theorem}

The non-zero components of $g^{\mathrm{F}}$ in the adapted coframe \smash{$\bigl(\tilde\th^0=\th, \tilde\th^1, \tilde\th^{2}, \tilde\th^{\ol1}, \tilde\th^{\ol2}, \tilde\th^3\bigr)$} are
\begin{equation}\label{metric-tensor}
g^{\mathrm{F}}_{1\ol2}=g^{\mathrm{F}}_{\ol2 1}=g^{\mathrm{F}}_{2\ol1}=g^{\mathrm{F}}_{\ol1 2}=g^{\mathrm{F}}_{03}=g^{\mathrm{F}}_{30}=1.
\end{equation}

\subsection[The connection form of g\^{}F]{The connection form of $\boldsymbol{g^{\mathrm{F}}}$}
The connection 1-forms $\omega_a{}^b$ of the Fefferman metric are computed as follows.

\begin{Proposition}\label{connection}
In the coframe \smash{$\bigl(\tilde\theta^0=\theta, \tilde\theta^1, \tilde\theta^{2}, \tilde\theta^{\ol1}, \tilde\theta^{\ol2}, \tilde\theta^3\bigr)$}, the connection forms of the Levi-Civita connection of $g^{\rm F}$ are given by
\begin{gather*}
\begin{alignedat}{7}
&\omega_0{}^0=0,\qquad&& \omega_1{}^0=\frac{{\rm i}}{2}\tilde\theta^{\ol2},\qquad&& \omega_2{}^0=\frac{{\rm i}}{2}\tilde\theta^{\ol1},\qquad&& \omega_3{}^0=0,& \\
&\omega_0{}^1=\frac{{\rm i}}{2}\tilde\theta^2,\qquad&& \omega_1{}^1=-\frac{{\rm i}}{2}\tilde\theta^3,\qquad&& \omega_2{}^1=-\frac{{\rm i}}{2}\theta,\qquad&& \omega_{\ol1}{}^1=0,\qquad && \omega_{\ol2}{}^1=0,\qquad&& \omega_3{}^1=\frac{{\rm i}}{2}\tilde\theta^1,&
\end{alignedat}
\\
\omega_0{}^2={\rm i}P(\zeta, \eta)\theta+{\rm i}P\bigl(\zeta, \ol\zeta\bigr)\tilde\theta^1+{\rm i}P(\zeta, \zeta)\tilde\theta^{\ol1},\\
\omega_1{}^2=-{\rm i}P(\eta, \eta)\theta-{\rm i}P\bigl(\ol\zeta, \eta\bigr)\tilde\theta^1-{\rm i}P(\zeta, \eta)\tilde\theta^{\ol1}, \\
\omega_2{}^2=-\frac{{\rm i}}{2}\tilde\theta^3, \qquad
\omega_{\ol1}{}^2=0, \qquad \omega_{\ol2}{}^2=0, \qquad \omega_3{}^2=\frac{{\rm i}}{2}\tilde\theta^2, \\
\omega_0{}^3=0, \qquad \omega_1{}^3={\rm i}P\bigl(\ol\zeta, \eta\bigr)\theta+{\rm i}P\bigl(\ol\zeta, \ol\zeta\bigr)\tilde\theta^1+{\rm i}P\bigl(\zeta, \ol\zeta\bigr)\tilde\theta^{\ol1}, \qquad \omega_2{}^3=\frac{{\rm i}}{2}\tilde\theta^{\ol2}, \qquad
\omega_3{}^3=0
\end{gather*}
and the complex conjugates of these forms.
\end{Proposition}
\begin{proof}
It suffices to check that these forms satisfy the structure equations
\[
g^{\mathrm{F}}_{bc}\omega_a{}^c+g^{\mathrm{F}}_{ac}\omega_b{}^c=dg^{\mathrm{F}}_{ab}, \qquad
{\rm d}\tilde\theta^a=\tilde\theta^b\wedge\omega_b{}^a.
\]
Since the metric tensor is given by \eqref{metric-tensor}, the first equation is equivalent to
\begin{alignat}{5}
&\omega_0{}^3=0,\qquad&& \omega_0{}^{\ol2}+\omega_1{}^3=0,\qquad&& \omega_0{}^{\ol1}+\omega_2{}^3=0,\qquad&& \omega_0{}^0+\omega_3{}^3=0, &\nonumber\\
&\omega_1{}^{\ol2}=0,\qquad&& \omega_1{}^{\ol1}+\omega_2{}^{\ol2}=0,\qquad&& \omega_1{}^{2}+\omega_{\ol1}{}^{\ol2}=0,\qquad&& \omega_1{}^1+\omega_{\ol2}{}^{\ol2}=0,& \nonumber\\
&\omega_1{}^0+\omega_3{}^{\ol2}=0,\qquad&& \omega_2{}^{\ol1}=0,\qquad&& \omega_2{}^{1}+\omega_{\ol2}{}^{\ol1}=0,\qquad&& \omega_2{}^0+\omega_3{}^{\ol1}=0,& \nonumber\\
&\omega_3{}^0=0&\label{metric-connection}
\end{alignat}
and these are satisfied. We can also verify the equation
\[
{\rm d}\tilde\theta^a=\theta\wedge\omega_0{}^a+\tilde\theta^1\wedge\omega_1{}^a+\tilde\theta^2\wedge\omega_2{}^a+\tilde\theta^{\ol1}\wedge\omega_{\ol1}{}^a+\tilde\theta^{\ol2}\wedge\omega_{\ol2}{}^a+\tilde\theta^3\wedge\omega_3{}^a
\]
for $a=0, 1, 2, 3$ by using \eqref{diff-coframe}.
\end{proof}

\subsection[The curvature form of g\^{}F]{The curvature form of $\boldsymbol{g^{\mathrm{F}}}$}
It is straightforward to compute the curvature form
\[
\Omega_a{}^b={\rm d}\omega_a{}^b-\omega_a{}^c\wedge\omega_c{}^b
=\frac{1}{2}R^{\mathrm{F}}_{cd}{}^b{}_a \tilde\theta^c\wedge\tilde\theta^d
\]
by using Proposition~\ref{connection} and the equations \eqref{d-eta} and \eqref{diff-coframe}.
Since $\Omega_a{}^b$ satisfies the same symmetries as in \eqref{metric-connection},
we present only
$\Omega_0{}^0$, $\Omega_1{}^0$, $\Omega_2{}^0$, $\Omega_0{}^{1}$, $\Omega_2{}^1$, $\Omega_0{}^2$, $\Omega_2{}^2$, $\Omega_1{}^2$, $\Omega_2{}^{\ol2}$.

\begin{Proposition}
The curvature forms of $g^{\mathrm{F}}$ are given by
\begin{gather*}
\Omega_0{}^0=\frac{1}{2}P\bigl(\ol\zeta, \eta\bigr)\theta\wedge\tilde\theta^{1}+\frac{1}{2}P(\zeta, \eta)\theta\wedge\tilde\theta^{\ol1}, \\
\Omega_1{}^0=\frac{1}{2}P\bigl(\ol\zeta, \ol\zeta\bigr)\theta\wedge\tilde\theta^1
 +\frac{1}{2}P\bigl(\zeta, \ol\zeta\bigr)\theta\wedge\tilde\theta^{\ol1}-\frac{1}{4}\tilde\theta^{\ol2}\wedge\tilde\theta^3, \\
 \Omega_2{}^0=\frac{1}{4}\theta\wedge\tilde\theta^{\ol2}-\frac{1}{4}\tilde\theta^{\ol1}\wedge\tilde\theta^3, \\
\Omega_0{}^1=-\frac{1}{2}P(\eta, \eta)\theta\wedge\tilde\theta^1+\frac{1}{2}P(\zeta, \eta)\tilde\theta^1\wedge\tilde\theta^{\ol1}+\frac{1}{4}\tilde\theta^2\wedge\tilde\theta^3, \\
\Omega_2{}^1=\frac{1}{4}\tilde\theta^1\wedge\tilde\theta^{\ol2}+\frac{1}{4}\tilde\theta^2\wedge\tilde\theta^{\ol1}, \\
\Omega_0{}^2=(*C)\bigl(\zeta, \ol\zeta\bigr)\theta\wedge\tilde\theta^1-\frac{1}{2}P(\eta, \eta)\theta\wedge\tilde\theta^2
 -(*C)(\zeta, \zeta)\theta\wedge\tilde\theta^{\ol1}+\frac{1}{2}P(\zeta, \eta)\theta\wedge\tilde\theta^3 \\
\hphantom{\Omega_0{}^2=}{}\hspace{0.2mm}
 -\frac{1}{2}P\bigl(\ol\zeta, \eta\bigr)\tilde\theta^1\wedge\tilde\theta^2
+(*C)(\zeta, \eta)\tilde\theta^1\wedge\tilde\theta^{\ol1}
+\frac{1}{2}P\bigl(\zeta, \ol\zeta\bigr) \tilde\theta^1\wedge\tilde\theta^3
+\frac{1}{2}P(\zeta, \eta)\tilde\theta^2\wedge\tilde\theta^{\ol1}
 \\
\hphantom{\Omega_0{}^2=}{}\hspace{0.2mm}
 +\frac{1}{2}P(\zeta, \zeta)\tilde\theta^{\ol1}\wedge\tilde\theta^3, \\
\Omega_2{}^2=\frac{1}{2}P(\zeta, \eta)\theta\wedge\tilde\theta^{\ol1}+\frac{1}{2}P\bigl(\zeta, \ol\zeta\bigr)\tilde\theta^1\wedge\tilde\theta^{\ol1}+\frac{1}{4}\tilde\theta^2\wedge\tilde\theta^{\ol2}, \\
\Omega_1{}^2=-(*C)\bigl(\ol\zeta, \eta\bigr)\theta\wedge\tilde\theta^1 -\frac{1}{2}P\bigl(\ol\zeta, \eta\bigr)\theta\wedge\tilde\theta^2
+(*C)(\zeta, \eta)\theta\wedge\tilde\theta^{\ol1}
+\frac{1}{2}P(\zeta, \eta)\theta\wedge\tilde\theta^{\ol2}
 \\
\hphantom{\Omega_1{}^2=}{}\hspace{0.2mm}
 -\frac{1}{2}P(\ol\zeta ,\ol\zeta)\tilde\theta^1\wedge\tilde\theta^2
-(*C)(\eta, \eta)\tilde\theta^1\wedge\tilde\theta^{\ol1}
+\frac{1}{2}
P\bigl(\zeta, \ol\zeta\bigr)\tilde\theta^1\wedge\tilde\theta^{\ol2}
+\frac{1}{2}P\bigl(\zeta, \ol\zeta\bigr)\tilde\theta^2\wedge\tilde\theta^{\ol1} \\
\hphantom{\Omega_1{}^2=}{}\hspace{0.2mm}
+\frac{1}{2}P(\zeta, \zeta)\tilde\theta^{\ol1}\wedge\tilde\theta^{\ol2}, \\
\Omega_2{}^{\ol2}=\frac{1}{2}P\bigl(\ol\zeta, \eta\bigr)\theta\wedge\tilde\theta^{\ol1}+\frac{1}{2}P\bigl(\ol\zeta, \ol\zeta\bigr)\tilde\theta^1\wedge\tilde\theta^{\ol1}.
\end{gather*}
The other components are computed by the symmetry and the reality of the curvature form.
\end{Proposition}
It follows that the non-zero components of the curvature tensor are
\begin{alignat*}{4}
&R^{\mathrm{F}}_{010\ol2}=\frac{1}{2}P(\eta, \eta),
\qquad&&
R^{\mathrm{F}}_{0\ol21\ol1}=-\frac{1}{2}P(\zeta, \eta),
\qquad&&
R^{\mathrm{F}}_{0\ol223}=-\frac{1}{4},& \\
&R^{\mathrm{F}}_{010\ol1}=-(*C)\bigl(\zeta, \ol\zeta\bigr),
\qquad&&
R^{\mathrm{F}}_{0\ol10\ol1}=(*C)(\zeta, \zeta),
\qquad&&
R^{\mathrm{F}}_{0\ol103}=-\frac{1}{2}P(\zeta, \eta),& \\
&R^{\mathrm{F}}_{0\ol112}=\frac{1}{2}P\bigl(\ol\zeta, \eta\bigr),
\qquad&&
R^{\mathrm{F}}_{0\ol12\ol1}=-\frac{1}{2}P(\zeta, \eta),
\qquad&&
R^{\mathrm{F}}_{0\ol11\ol1}=-(*C)(\zeta, \eta),& \\
&R^{\mathrm{F}}_{0\ol113}=-\frac{1}{2}P\bigl(\zeta, \ol\zeta\bigr),
\qquad&&
R^{\mathrm{F}}_{0\ol1\ol13}=-\frac{1}{2}P(\zeta, \zeta),
\qquad&&
R^{\mathrm{F}}_{121\ol1}=\frac{1}{2}P\bigl(\ol\zeta, \ol\zeta\bigr),& \\
&R^{\mathrm{F}}_{2\ol12\ol2}=-\frac{1}{4},
\qquad&&
R^{\mathrm{F}}_{1\ol12\ol1}=-\frac{1}{2}P\bigl(\zeta, \ol\zeta\bigr),
\qquad&&
R^{\mathrm{F}}_{23\ol13}=\frac{1}{4},& \\
&R^{\mathrm{F}}_{1\ol11\ol1}=(*C)(\eta, \eta)&
\end{alignat*}
and the components which can be obtained from these by the symmetry
$R^{\mathrm{F}}_{abcd}=R^{\mathrm{F}}_{[ab][cd]}=R^{\mathrm{F}}_{cdab}$ and reality of the curvature tensor.

The non-zero components of the Ricci tensor
\[
R^{\mathrm{F}}_{ab}=R^{\mathrm{F}}_{1a\ol2b}+R^{\mathrm{F}}_{\ol2a1b}+R^{\mathrm{F}}_{2a\ol1b}+R^{\mathrm{F}}_{\ol1a2b}+R^{\mathrm{F}}_{0a3b}+R^{\mathrm{F}}_{3a0b}
\]
are given by
\begin{alignat*}{4}
&R^{\mathrm{F}}_{00}=2P(\eta, \eta), \qquad&&
R^{\mathrm{F}}_{01}=2P\bigl(\ol\zeta, \eta\bigr), \qquad&&
R^{\mathrm{F}}_{2\ol2}=1,& \\
&R^{\mathrm{F}}_{11}=2P\bigl(\ol\zeta, \ol\zeta\bigr),\qquad&&
R^{\mathrm{F}}_{1\ol1}=2P\bigl(\zeta, \ol\zeta\bigr),\qquad&&
R^{\mathrm{F}}_{33}=1.&
\end{alignat*}
It follows that the scalar curvature vanishes,
\[
R^{\mathrm{F}}=2\bigl(R^{\mathrm{F}}_{1\ol2}+R^{\mathrm{F}}_{2\ol1}+R^{\mathrm{F}}_{03}\bigr)=0.
\]
The Schouten tensor is equal to $P^{\mathrm{F}}_{ab}=\frac{1}{4}R^{\mathrm{F}}_{ab}$, so the non-zero components of the Weyl curvature
\[
W^{\mathrm{F}}_{abcd}=R^{\mathrm{F}}_{abcd}-\frac{1}{4}
\bigl(R^{\mathrm{F}}_{ac}g^{\mathrm{F}}_{bd}-R^{\mathrm{F}}_{bc}g^{\mathrm{F}}_{ad}+R^{\mathrm{F}}_{bd}g^{\mathrm{F}}_{ac}-R^{\mathrm{F}}_{ad}g^{\mathrm{F}}_{bc}\bigr)
\]
are given by
\begin{alignat*}{3}
&W^{\mathrm{F}}_{010\ol1}=-(*C)\bigl(\zeta, \ol\zeta\bigr),\qquad&&
W^{\mathrm{F}}_{0\ol10\ol1}=(*C)(\zeta, \zeta),& \\
&W^{\mathrm{F}}_{0\ol11\ol1}=-(*C)(\zeta, \eta),\qquad&&
W^{\mathrm{F}}_{1\ol11\ol1}=(*C)(\eta, \eta)&
\end{alignat*}
and the components which can be obtained from these by the symmetry and reality of the Weyl tensor.
Since $\bigl(\zeta, \ol\zeta, \eta\bigr)$ gives a basis of $\mathbb{C}T\Sigma$ at the base point, we obtain the following theorem.

\begin{Theorem}[{\cite[Theorem 6.7]{Low}} and \cite{SY}]\label{flat}
The Fefferman space $\bigl(\wh{\mathcal{C}}, g^{\mathrm{F}}\bigr)$ is conformally flat if and only if $(\Sigma, [g])$ is conformally flat.
\end{Theorem}

This also implies that the twistor CR manifold $M$ is locally CR equivalent to the hyperquadric if and only if $(\Sigma, [g])$ is conformally flat.

\subsection[Null geodesics of g\^{}F]{Null geodesics of $\boldsymbol{g^{\mathrm{F}}}$}
To examine the projections of chains and null chains to $\Sigma$, we
compute the null geodesic equation for the Fefferman metric.
We fix a representative metric $g\in[g]$, and let \smash{$\gamma(t)=\bigl(x^i(t), \zeta_i(t)\bigr)$} be a curve on $\wh{\mathcal{C}}$ such that the projection \smash{$x(t)\!=\!\bigl(x^i(t)\bigr)$} is a regular curve on $\Sigma$. In the adapted coframe
$\smash{\bigl(\tilde\theta^0\!=\!\theta, \tilde\theta^1, \tilde\theta^2, \tilde\theta^{\ol1}, \tilde\theta^{\ol2}, \tilde\theta^3\bigr)}$, the components \smash{$\dot\gamma^a:=\tilde\th^a(\dot\gamma)$} are given by
\begin{align*}
&\dot\gamma^0=\eta_i \dot x^i=g(\eta, \dot x), \\
&\dot\gamma^1=\zeta_i \dot x^i=g(\zeta, \dot x), \\
&\dot\gamma^2={\rm i}\eta^j\bigl(\dot\zeta_j-\Gamma_{kj}{}^l\zeta_l \dot x^k\bigr)
={\rm i}g(\eta, \nabla_{\dot x}\zeta), \\
&\dot\gamma^3=-\frac{{\rm i}}{2}\bigl(
\ol{\zeta^j}\bigl(\dot\zeta_j-\Gamma_{kj}{}^l\zeta_l \dot x^k\bigr)
-\zeta^j\bigl(\ol{\dot\zeta_j}-\Gamma_{kj}{}^l\ol{\zeta_l} \dot x^k\bigr)
\bigr)=\operatorname{Im} g\bigl(\ol\zeta, \nabla_{\dot x}\zeta\bigr)
=-{\rm i} g\bigl(\ol\zeta, \nabla_{\dot x}\zeta\bigr),
\end{align*}
where $\nabla$ is the Levi-Civita connection of $g$.
The nullity of $\gamma$ is written as
\begin{equation}\label{nullity}
g^{\rm F}(\dot\gamma, \dot\gamma)=2\bigl(\dot\gamma^1\dot\gamma^{\ol2}+\dot\gamma^2\dot\gamma^{\ol1}\bigr)+2\dot\gamma^0\dot\gamma^3=0.
\end{equation}
Let $\nabla^{\mathrm{F}}$ be the Levi-Civita connection of $g^{\mathrm{F}}$. By Proposition~\ref{connection}, we can compute the acceleration
\[
\nabla^{\mathrm{F}}_{\dot\gamma}\dot\gamma^a
=\ddot\gamma^a+\omega_0{}^a(\dot\gamma)\dot\gamma^0
+\omega_1{}^a(\dot\gamma)\dot\gamma^1+\omega_2{}^a(\dot\gamma)\dot\gamma^2+\omega_{\ol1}{}^a(\dot\gamma)\dot\gamma^{\ol1}+\omega_{\ol2}{}^a(\dot\gamma)\dot\gamma^{\ol2}+\omega_3{}^a(\dot\gamma)\dot\gamma^3
\]
as follows:
\begin{gather*}
\nabla^{\mathrm{F}}_{\dot\gamma}\dot\gamma^0=\ddot\gamma^0, \\
\nabla^{\mathrm{F}}_{\dot\gamma}\dot\gamma^1=\ddot\gamma^1, \\
\nabla^{\mathrm{F}}_{\dot\gamma}\dot\gamma^2
= \ddot\gamma^2+{\rm i}P(\zeta, \eta)\bigl(\dot\gamma^0\bigr)^2+{\rm i}\bigl(P\bigl(\zeta, \ol\zeta\bigr)-P(\eta, \eta)\bigr)\dot\gamma^0\dot\gamma^1+{\rm i}P(\zeta, \zeta)\dot\gamma^0\dot\gamma^{\ol1} \\
\hphantom{\nabla^{\mathrm{F}}_{\dot\gamma}\dot\gamma^2=}{}
-{\rm i}P\bigl(\ol\zeta, \eta\bigr)\bigl(\dot\gamma^1\bigr)^2-{\rm i}P(\zeta, \eta)|\dot\gamma^1|^2, \\
\nabla^{\mathrm{F}}_{\dot\gamma}\dot\gamma^3
=\ddot\gamma^3+{\rm i}P\bigl(\ol\zeta, \eta\bigr)\dot\gamma^0\dot\gamma^1-{\rm i}P(\zeta, \eta)\dot\gamma^0\dot\gamma^{\ol1}+{\rm i}P\bigl(\ol\zeta, \ol\zeta\bigr)\bigl(\dot\gamma^1\bigr)^2-{\rm i}P(\zeta, \zeta)\bigl(\dot\gamma^{\ol1}\bigr)^2.
\end{gather*}
Hence, $\gamma$ is a geodesic of $g^{\mathrm{F}}$ if and only if
\begin{equation*}
\dot\gamma^0=\mathrm{const}, \qquad \dot\gamma^1=\mathrm{const},
\end{equation*}
and
\begin{gather}
 \ddot\gamma^2+{\rm i}P(\zeta, \eta)\bigl(\dot\gamma^0\bigr)^2+{\rm i}\bigl(P\bigl(\zeta, \ol\zeta\bigr)-P(\eta, \eta)\bigr)\dot\gamma^0\dot\gamma^1+{\rm i}P(\zeta, \zeta) \dot\gamma^0\dot\gamma^{\ol1} \nonumber \\
 \qquad{} -{\rm i}P\bigl(\ol\zeta, \eta\bigr)\bigl(\dot\gamma^1\bigr)^2-{\rm i}P(\zeta, \eta)|\dot\gamma^1|^2=0,\label{geodesic2}
\\
\ddot\gamma^3+{\rm i}P\bigl(\ol\zeta, \eta\bigr)\dot\gamma^0\dot\gamma^1-{\rm i}P(\zeta, \eta)\dot\gamma^0\dot\gamma^{\ol1}+{\rm i}P\bigl(\ol\zeta, \ol\zeta\bigr)\bigl(\dot\gamma^1\bigr)^2-{\rm i}P(\zeta, \zeta)\bigl(\dot\gamma^{\ol1}\bigr)^2=0.\label{geodesic3}
\end{gather}

We will derive the equation satisfied by $x(t)$ from the null geodesic equation for $\gamma(t)$. We note that for any $V, W\in \mathbb{C}T_x\Sigma$, the identities
\begin{align*}
&V=g(V, \eta)\eta+g\bigl(V, \ol\zeta\bigr)\zeta+g(V, \zeta)\ol\zeta, \\
&g(V, W)=g(V, \eta)g(W, \eta)+g(V, \zeta)g\bigl(W, \ol\zeta\bigr)+g\bigl(V, \ol\zeta\bigr)g(W, \zeta)
\end{align*}
hold on $\hat\pi^{-1}(x)\subset\wh{\mathcal{C}}$. In particular, we have
\begin{equation}\label{x-zeta-eta}
\dot x=\dot\gamma^0\eta+\dot\gamma^{\ol1}\zeta+\dot\gamma^1\ol\zeta, \qquad |\dot x|^2=\bigl(\dot\gamma^0\bigr)^2+2\bigl|\dot\gamma^1\bigr|^2.
\end{equation}
\begin{Lemma}
For any curve $\gamma(t)=(x(t), \zeta(t))$ on $\wh{\mathcal{C}}$, it holds that
\begin{alignat}{3}
&g\bigl(\nabla_{\dot x}\nabla_{\dot x}\zeta, \ol\zeta\bigr) ={\rm i}\ddot\gamma^3-\bigl|\dot\gamma^2\bigr|^2-\bigl(\dot\gamma^3\bigr)^2, \qquad&&
g(\nabla_{\dot x}\nabla_{\dot x}\zeta, \zeta) =\bigl(\dot\gamma^2\bigr)^2,& \nonumber \\
&g(\nabla_{\dot x}\nabla_{\dot x}\zeta, \eta) =-{\rm i}\ddot\gamma^2+\dot\gamma^2\dot\gamma^3, \qquad&&
g(\nabla_{\dot x}\nabla_{\dot x}\eta, \zeta) ={\rm i}\ddot\gamma^2+\dot\gamma^2\dot\gamma^3, &\nonumber \\
&g(\nabla_{\dot x}\nabla_{\dot x}\eta, \eta) =-2\bigl|\dot\gamma^2\bigr|^2.&\label{nabla-zeta-eta}
\end{alignat}
\end{Lemma}
\begin{proof}
We only prove the identity for $g(\nabla_{\dot x}\nabla_{\dot x}\eta, \zeta)$ since the other cases are similar. Using the identities noted above, we have
\begin{align*}
g(\nabla_{\dot x}\nabla_{\dot x}\eta, \zeta)
&=\frac{{\rm d}}{{\rm d}t}g(\nabla_{\dot x}\eta, \zeta)-g(\nabla_{\dot x}\eta, \nabla_{\dot x}\zeta) =-\frac{{\rm d}}{{\rm d}t}g(\eta, \nabla_{\dot x}\zeta)-g(\nabla_{\dot x}\eta, \zeta)g\bigl(\nabla_{\dot x}\zeta, \ol\zeta\bigr) \\
&=-\frac{{\rm d}}{{\rm d}t}g(\eta, \nabla_{\dot x}\zeta)+g(\eta, \nabla_{\dot x}\zeta)g\bigl(\nabla_{\dot x}\zeta, \ol\zeta\bigr)
={\rm i}\ddot\gamma^2+\dot\gamma^2\dot\gamma^3.
\tag*{\qed}
\end{align*}
\renewcommand{\qed}{}
\end{proof}

Now we are ready to prove the following.

\begin{Theorem}\label{null-geod-proj}
Let $g\in[g]$ be a representative metric and let $\gamma(t)=(x(t), \zeta(t))$ be a null geodesic of $g^{\mathrm{F}}$ on $\wh{\mathcal{C}}$ such that the projection $x(t)$ is a regular curve on $\Sigma$. Then, $|\dot x|_g$ is constant, and $x(t)$ satisfies
\begin{equation}\label{eq-x}
\nabla_{\dot x}\nabla_{\dot x}\dot x-|\dot x|^2P(\dot x)=
-\bigl(2\bigl|\dot\gamma^2\bigr|^2+\bigl(\dot\gamma^3\bigr)^2+P(\dot x, \dot x)\bigr)\dot x.
\end{equation}
In particular, $x(t)$ is a conformal geodesic of $(\Sigma, [g])$.
\end{Theorem}
\begin{proof}
We prove that the inner products of both sides of \eqref{eq-x} with $\zeta$, $\eta$ coincide.

Since $\gamma$ is a geodesic, $\dot\gamma^0$ and $\dot\gamma^1$ are constant.
Hence, by \eqref{x-zeta-eta}, $|\dot x|$ is constant and we have
\begin{align*}
\nabla_{\dot x}\nabla_{\dot x}\dot x-|\dot x|^2P(\dot x)
=&\dot\gamma^0\nabla_{\dot x}\nabla_{\dot x}\eta+
\dot\gamma^{\ol1}\nabla_{\dot x}\nabla_{\dot x}\zeta+\dot\gamma^1\nabla_{\dot x}\nabla_{\dot x}\ol\zeta \\
&{}{-}\,\bigl(\bigl(\dot\gamma^0\bigr)^2+2\bigl|\dot\gamma^1\bigr|^2\bigr)\bigl(\dot\gamma^0P(\eta)+\dot\gamma^{\ol1}P(\zeta)+\dot\gamma^1P(\ol\zeta)\bigr).
\end{align*}
By using \eqref{nabla-zeta-eta}, we compute as
\begin{align*}
g\bigl(\nabla_{\dot x}\nabla_{\dot x}\dot x-|\dot x|^2P(\dot x) ,\zeta\bigr)
={}&\dot\gamma^0g(\nabla_{\dot x}\nabla_{\dot x}\eta, \zeta)+
\dot\gamma^{\ol1}g(\nabla_{\dot x}\nabla_{\dot x}\zeta, \zeta)+\dot\gamma^1g\bigl(\nabla_{\dot x}\nabla_{\dot x}\ol\zeta, \zeta\bigr) \\
&{}{-}\,\bigl(\bigl(\dot\gamma^0\bigr)^2+2\bigl|\dot\gamma^1\bigr|^2\bigr)\bigl(\dot\gamma^0P(\eta, \zeta)+\dot\gamma^{\ol1}P(\zeta, \zeta)+\dot\gamma^1P\bigl(\ol\zeta, \zeta\bigr)\bigr) \\
={}&{\rm i}\dot\gamma^0\ddot\gamma^2+\dot\gamma^0\dot\gamma^2\dot\gamma^3+\dot\gamma^{\ol1}\bigl(\dot\gamma^2\bigr)^2
-{\rm i}\dot\gamma^1\ddot\gamma^3-\dot\gamma^1\bigl|\dot\gamma^2\bigr|^2-\dot\gamma^1\bigl(\dot\gamma^3\bigr)^2 \\
&{}{-}\bigl(\bigl(\dot\gamma^0\bigr)^2+2|\dot\gamma^1|^2\bigr)\bigl(\dot\gamma^0P(\zeta, \eta)+\dot\gamma^{\ol1}P(\zeta, \zeta)+\dot\gamma^1P\bigl(\zeta, \ol\zeta\bigr)\bigr).
\end{align*}
By the nullity \eqref{nullity}, we have
\[
\dot\gamma^0\dot\gamma^2\dot\gamma^3+\dot\gamma^{\ol1}\bigl(\dot\gamma^2\bigr)^2 -\dot\gamma^1\bigl|\dot\gamma^2\bigr|^2=-2\dot\gamma^1\bigl|\dot\gamma^2\bigr|^2.
\]
Moreover, the geodesic equations \eqref{geodesic2} and \eqref{geodesic3} give
\begin{gather*}
{\rm i}\dot\gamma^0\ddot\gamma^2=P(\zeta, \eta)\bigl(\dot\gamma^0\bigr)^3+\bigl(P\bigl(\zeta, \ol\zeta\bigr)-P(\eta, \eta)\bigr)\bigl(\dot\gamma^0\bigr)^2\dot\gamma^1+P(\zeta, \zeta) \bigl(\dot\gamma^0\bigr)^2\dot\gamma^{\ol1} \\
\hphantom{{\rm i}\dot\gamma^0\ddot\gamma^2=}{}
-P\bigl(\ol\zeta, \eta\bigr)\dot\gamma^0\bigl(\dot\gamma^1\bigr)^2-P(\zeta, \eta)\dot\gamma^0\bigl|\dot\gamma^1\bigr|^2, \\
-{\rm i}\dot\gamma^1\ddot\gamma^3=-P\bigl(\ol\zeta, \eta\bigr)\dot\gamma^0\bigl(\dot\gamma^1\bigr)^2+P(\zeta, \eta)\dot\gamma^0\bigl|\dot\gamma^{1}\bigr|^2
-P\bigl(\ol\zeta, \ol\zeta\bigr)\bigl(\dot\gamma^1\bigr)^3+P(\zeta, \zeta)\dot\gamma^{\ol1}\bigl|\dot\gamma^{1}\bigr|^2.
\end{gather*}
Substituting these expressions, we obtain
\begin{align*}
g\bigl(\nabla_{\dot x}\nabla_{\dot x}\dot x-|\dot x|^2P(\dot x) ,\zeta\bigr)
&=-2\dot\gamma^1\bigl|\dot\gamma^2\bigr|^2-\dot\gamma^1\bigl(\dot\gamma^3\bigr)^2-\dot\gamma^1P(\dot x, \dot x) \\
&=-\bigl(2\bigl|\dot\gamma^2\bigr|^2+\bigl(\dot\gamma^3\bigr)^2+P(\dot x, \dot x)\bigr)g(\dot x, \zeta).
\end{align*}

Similarly, we compute as
\begin{align*}
g\bigl(\nabla_{\dot x}\nabla_{\dot x}\dot x-|\dot x|^2P(\dot x) ,\eta\bigr)
={}&\dot\gamma^0g(\nabla_{\dot x}\nabla_{\dot x}\eta, \eta)+
\dot\gamma^{\ol1}g(\nabla_{\dot x}\nabla_{\dot x}\zeta, \eta)+\dot\gamma^1g\bigl(\nabla_{\dot x}\nabla_{\dot x}\ol\zeta, \eta\bigr) \\
&{}{-}\bigl(\bigl(\dot\gamma^0\bigr)^2+2\bigl|\dot\gamma^1\bigr|^2\bigr)\bigl(\dot\gamma^0P(\eta, \eta)+\dot\gamma^{\ol1}P(\zeta, \eta)+\dot\gamma^1P\bigl(\ol\zeta, \eta\bigr)\bigr) \\
={}&{-}2\dot\gamma^0\bigl|\dot\gamma^2\bigr|^2-{\rm i}\dot\gamma^{\ol1}\ddot\gamma^2
+\dot\gamma^{\ol1}\dot\gamma^2\dot\gamma^3+{\rm i}\dot\gamma^1\ddot\gamma^{\ol2}+\dot\gamma^1\dot\gamma^{\ol2}\dot\gamma^3 \\
&{}{-}\bigl(\bigl(\dot\gamma^0\bigr)^2+2\bigl|\dot\gamma^1\bigr|^2\bigr)\bigl(\dot\gamma^0P(\eta, \eta)+\dot\gamma^{\ol1}P(\zeta, \eta)+\dot\gamma^1P\bigl(\ol\zeta, \eta\bigr)\bigr).
\end{align*}
By the nullity and the equation \eqref{geodesic2}, we have
\begin{align*}
\dot\gamma^1\dot\gamma^{\ol2}\dot\gamma^3+\dot\gamma^{\ol1}\dot\gamma^2\dot\gamma^3 =-\dot\gamma^0\bigl(\dot\gamma^3\bigr)^2
\end{align*}
and
\begin{align*}
-{\rm i}\dot\gamma^{\ol1}\ddot\gamma^2+{\rm i}\dot\gamma^1\ddot\gamma^{\ol2}
={}&{-}P(\zeta, \eta)\bigl(\dot\gamma^0\bigr)^2\dot\gamma^{\ol1} -P\bigl(\ol\zeta, \eta\bigr)\bigl(\dot\gamma^0\bigr)^2\dot\gamma^1-2\bigl(P\bigl(\zeta, \ol\zeta\bigr)-P(\eta, \eta)\bigr)
\dot\gamma^0\bigl|\dot\gamma^1\bigr|^2\\
&{}{-}\, P(\zeta, \zeta) \dot\gamma^0\bigl(\dot\gamma^{\ol1}\bigr)^2\!-P\bigl(\ol\zeta, \ol\zeta\bigr) \dot\gamma^0\bigl(\dot\gamma^{1}\bigr)^2+2P\bigl(\ol\zeta, \eta\bigr)\dot\gamma^1\bigl|\dot\gamma^1\bigr|^2
 + 2P(\zeta, \eta)\dot\gamma^{\ol1}\bigl|\dot\gamma^1\bigr|^2.
\end{align*}
Substituting these expressions, we obtain
\begin{align*}
g\bigl(\nabla_{\dot x}\nabla_{\dot x}\dot x-|\dot x|^2P(\dot x) ,\eta\bigr)
&=-2\dot\gamma^0\bigl|\dot\gamma^2\bigr|^2-\dot\gamma^0\bigl(\dot\gamma^3\bigr)^2-\dot\gamma^0P(\dot x, \dot x) \\
&=-\bigl(2\bigl|\dot\gamma^2\bigr|^2+\bigl(\dot\gamma^3\bigr)^2+P(\dot x, \dot x)\bigr)g(\dot x, \eta).
\end{align*}
Thus, $x(t)$ satisfies \eqref{eq-x}. In particular, it satisfies the conformal geodesic equation \eqref{conf-geod-const-speed} for curves with constant speed.
\end{proof}

By definition, a (null) chain $(x(t), [\zeta(t)])$ is the projection of a non-vertical null geodesic $\gamma(t)=(x(t), \zeta(t))$ on $\wh{\mathcal{C}}$ to $M$. If the constants $\dot\gamma^0=g(\eta, \dot x)$,
$\dot\gamma^1=g(\zeta, \dot x)$ are both equal to~0, then we have $\dot x=0$ and the projection $x(t)$ is a~constant curve on $\Sigma$. Otherwise, it is a regular curve and hence a conformal geodesic on $\Sigma$ by the above theorem. Thus we obtain Theorem~\ref{chain-projection}\,(1).

As an application, we have the following.

\begin{Theorem}[cf.\ \cite{SY}]
Let $M_i$, $i=1, 2$ be the twistor CR manifolds over conformal $3$-manifolds $(\Sigma_i, [g_i])$, $i=1, 2$. If $\wh f\colon M_1\to M_2$ is a bundle isomorphism which is a CR equivalent map, then the underlying map $f\colon (\Sigma_1, [g_1])\to (\Sigma_2, [g_2])$ is a conformal isomorphism.
\end{Theorem}
\begin{proof}
Since $\wh f$ maps any (null) chain on $M_1$ to a (null) chain on $M_2$, $f$ maps conformal geodesics on $\Sigma_1$ to conformal geodesics on $\Sigma_2$. Hence, by~\cite[Theorem 2.3]{Kuo}, $f$ is a conformal isomorphism.
\end{proof}

\subsection{Canonical lift of conformal geodesics}\label{lift}
Theorem~\ref{chain-projection}\,(2) asserts that we can lift a conformal geodesic $x(t)$ to a unique (null) chain on $M$ once we choose a lift of the endpoint $x(a)$; we prove this in Section~\ref{sphere-bundle} below.
Here, as a special case, we present a canonical lift of a conformal geodesic to a chain.

Given a regular curve $x(t)$ on $\Sigma$, we take a complex vector field $\zeta(t)$ along $x(t)$ such that $(\sqrt{2}\operatorname{Re}\zeta, \sqrt{2}\operatorname{Im}\zeta)$ gives an oriented orthonormal basis of $\dot x^\perp$ with respect to a metric ${g\in[g]}$. Such a $\zeta(t)$ is determined by the conformal structure $[g]$ up to $\mathbb{C}^*$-action, so we have a canonical~lift
\[
\tilde x(t):=(x(t), [\zeta(t)])
\]
on $M$. Note that $\tilde x(t)$ is transverse to the contact distribution since
$\th(\dot{\tilde x})=|\dot x|^2\neq 0$.
\begin{Theorem} \label{canonical-lift}
The curve $\tilde x(t)$ is a chain if and only if $x(t)$ is a conformal geodesic. Moreover, for a conformal geodesic $x(t)$, $a\le t\le b$, the curve $\tilde x(t)$ is a unique chain through $(x(a), [\zeta(a)])$ which projects to $x(t)$.
\end{Theorem}
\begin{proof}
If $\tilde x(t)$ is a chain, then $x(t)$ is a conformal geodesic by Theorem~\ref{chain-projection}\,(1). Conversely, suppose $x(t)$ is a conformal geodesic. By the characterization of conformal geodesics given by Proposition~\ref{conf-geod-characterization}, we may take a representative metric $g\in[g]$ and a parametrization $t$ such that $\nabla_{\dot x}\dot x=P(\dot x)=0$ and $|\dot x|=1$. Since $\dot x$ is parallel, we can take
$\zeta(t)$ as a parallel vector field along~$x(t)$.
Then, $\gamma(t):=(x(t), \zeta(t))$ becomes a curve on $\wh{\mathcal{C}}$ with
\[
\eta=\dot x, \qquad \dot\gamma^0=1, \qquad \dot\gamma^1=\dot\gamma^2=\dot\gamma^3=0.
\]
In particular, it is a null curve. Noting that $P(\eta)=0$, we see that $\gamma(t)$ satisfies the geodesic equations \eqref{geodesic2} and \eqref{geodesic3}, and hence $\tilde x(t)=(x(t), [\zeta(t)])$ is a chain on $M$.

Next, we prove the uniqueness. Suppose that $c (t)$ is a chain through $(x(a), [\zeta(a)])$ which projects to a conformal geodesic $x(t)$, and let $\gamma_c(t)=(x(t), \zeta_c(t))$ be a null geodesic on~${\bigl(\wh{\mathcal{C}}, g^{\mathrm{F}}\bigr)}$, the Fefferman space determined by $g$ above, which projects to $c(t)$. Then, $\dot\gamma_c^0=g(\eta_c, \dot x)$ and $\dot\gamma_c^1=g(\zeta_c, \dot x)$ are constant by the geodesic equation. Since $\dot\gamma_c^0(a)=|\dot x(a)|^2=1$, $\dot\gamma_c^1(a)=0$, we obtain that $\eta_c(t)=\dot x(t)$ for any $t$ and hence $c(t)=\tilde x(t)$.
\end{proof}

\section[A variational principle for conformal geodesics in dimension three]{A variational principle for conformal geodesics\\ in dimension three}\label{variational-principle}

\subsection{Relations to the unit tangent sphere bundle and the frame bundle}\label{sphere-bundle}
We will give an alternative description of the previous constructions in terms of the frame bundle over $\Sigma$.
We fix a representative metric $g\in[g]$ and let
$\mathcal{P}$ be the $\mathrm{SO}(3)$-bundle of oriented orthonormal frames over $(\Sigma, g)$. Then, the map
\begin{equation}\label{isom-C-hat}
f\colon\ \wh{\mathcal{C}}\overset{\sim}{\longrightarrow} \mathcal{P}, \qquad
\zeta \longmapsto B=\bigl(\sqrt{2}\operatorname{Re}\zeta, \sqrt{2}\operatorname{Im}\zeta, \eta\bigr)
\end{equation}
gives a bundle isomorphism. The action of $\mathrm{U}(1)$ on $\wh{\mathcal{C}}$ corresponds to rotations of the first two basis vectors in $B$. Thus, the twistor CR manifold $M= \wh{\mathcal{C}}/\mathrm{U}(1)$ is isomorphic to the unit tangent sphere bundle of $(\Sigma, g)$ by the following map:
\begin{equation}\label{M-isom}
M \overset{\sim}{\longrightarrow} S\Sigma, \qquad
[\zeta]\longmapsto \eta.
\end{equation}



Let
\[
\phi=\begin{pmatrix}
\phi^1 \\ \phi^2 \\ \phi^3
\end{pmatrix}
\colon\ T\mathcal{P}\longrightarrow \mathbb{R}^3, \qquad \omega=\begin{pmatrix}
\hphantom{-}0 & \hphantom{-}\omega^3 & -\omega^2 \\
-\omega^3 & \hphantom{-}0 & \hphantom{-}\omega^1 \\
\hphantom{-}\omega^2 & -\omega^1 & \hphantom{-}0
\end{pmatrix}
\colon\ T\mathcal{P}\longrightarrow \mathfrak{so}(3)
\]
be the canonical $1$-form and the Levi-Civita connection form of $g$.
These forms are related to the adapted coframe \smash{$\bigl\{\theta, \tilde\theta^1, \tilde\theta^2, \tilde\theta^{\ol1}, \tilde\theta^{\ol2}, \tilde\theta^3\bigr\}$} on $\wh{\mathcal{C}}$ as follows.
\begin{Lemma}
We have
\[
f^*\phi=
\begin{pmatrix}
\sqrt{2}\operatorname{Re}\tilde\theta^1 \\
\sqrt{2}\operatorname{Im}\tilde\theta^1 \\
\theta
\end{pmatrix}, \qquad
f^*\omega=\begin{pmatrix}
0 & \tilde\theta^3 & -\sqrt{2}\operatorname{Im}\tilde\theta^2 \\
-\tilde\theta^3 & 0 & \hphantom{-}\sqrt{2}\operatorname{Re}\tilde\theta^2 \\
\sqrt{2}\operatorname{Im}\tilde\theta^2 & -\sqrt{2}\operatorname{Re}\tilde\theta^2& 0
\end{pmatrix}.
\]
\end{Lemma}
\begin{proof}
We fix a point $x_0\in\Sigma$ and take normal coordinates $(x^i)$ of $g$ around
$x_0$. Then, at any point $\zeta\in \wh{\mathcal{C}}_{x_0}$, we have
\[
f^*\phi={}^t \!B\, {\rm d}x =\begin{pmatrix}
\sqrt{2}\operatorname{Re}{}^t\zeta \\
\sqrt{2}\operatorname{Im}{}^t\zeta \\
{}^t\eta
\end{pmatrix}
\begin{pmatrix}
{\rm d}x^1 \\ {\rm d}x^2 \\ {\rm d}x^3
\end{pmatrix}
\]
and
\[
f^*\omega={}^t\!B\,{\rm d}B=
\begin{pmatrix}
\sqrt{2}\operatorname{Re}{}^t\zeta \\
\sqrt{2}\operatorname{Im}{}^t\zeta \\
{}^t\eta
\end{pmatrix}
\begin{pmatrix}
\sqrt{2}\operatorname{Re}{\rm d}\zeta & \sqrt{2}\operatorname{Im} {\rm d}\zeta & {\rm d}\eta
\end{pmatrix},
\]
where we write $\zeta$, $\eta$ as column vectors. Thus we obtain the desired formulas.
\end{proof}

We define the Fefferman metric $G^{\mathrm{F}}$ on $\mathcal{P}$ via the isomorphism \eqref{isom-C-hat}: $g^{\mathrm{F}}=f^* G^{\mathrm{F}}$. Then, the above lemma and Theorem~\ref{Fefferman-metric} prove the following.

\begin{Proposition}
The Fefferman metric $G^{\mathrm{F}}$ on $\mathcal{P}$ is given by
\[
G^{\mathrm{F}}=2\bigl( \phi^1\cdot\omega^1+ \phi^2\cdot\omega^2+ \phi^3\cdot\omega^3\bigr).
\]
\end{Proposition}

\begin{Remark}\quad
\begin{itemize}\itemsep=0pt
\item[(i)]
If we denote by $R_A$ the right action of $A\in\mathrm{SO}(3)$ on $\mathcal{P}$, we have
\[
R_A^*\begin{pmatrix}
\phi^1 \\ \phi^2 \\ \phi^3
\end{pmatrix}=A^{-1}\begin{pmatrix}
\phi^1 \\ \phi^2 \\ \phi^3
\end{pmatrix}, \qquad R_A^*\begin{pmatrix}
\omega^1 \\ \omega^2 \\ \omega^3
\end{pmatrix}=A^{-1}\begin{pmatrix}
\omega^1 \\ \omega^2 \\ \omega^3
\end{pmatrix}.
\]
Hence, $G^{\mathrm{F}}$ is a right $\mathrm{SO}(3)$-invariant metric: $R_A^*G^{\mathrm{F}}=G^{\mathrm{F}}$.

\item[(ii)] Since the conformal class of the Fefferman metric is invariantly defined from the CR structure, the conformal manifold $\bigl(\mathcal{P}, \bigl[G^{\mathrm{F}}\bigr]\bigr)$ should also be canonically associated to $(\Sigma, [g])$. This can be checked directly as follows.
Let $\wh g={\rm e}^{2\Upsilon}g$ be another representative metric and~$\wh{\mathcal{P}}$ be the oriented orthonormal frame bundle for $\wh g$.
Then, we have the canonical bundle isomorphism
\[
\d\colon\ \mathcal{P}\overset{\sim}{\longrightarrow}
\wh{\mathcal{P}}, \qquad B\longmapsto {\rm e}^{-\Upsilon}B.
\]
We denote by
$\wh\phi$, $\wh\omega$ the canonical 1-form and the connection form of $\wh g$ on $\wh{\mathcal{P}}$. Then, we have $\d^*\wh\phi={\rm e}^{\Upsilon}\phi$ and
the conformal transformation formula for the Levi-Civita connection gives
\begin{align*}
&\d^*\wh\omega^1=\omega^1+\Upsilon^3\phi^2-\Upsilon^2\phi^3, \\
&\d^*\wh\omega^2=\omega^2+\Upsilon^1\phi^3-\Upsilon^3\phi^1, \\
&\d^*\wh\omega^3=\omega^3+\Upsilon^2\phi^1-\Upsilon^1\phi^2,
\end{align*}
where we set $\Upsilon^i:=\phi^i(\mathrm{grad}_g \Upsilon)\in C^\infty(\mathcal{P})$, which is well-defined since $\phi$ is horizontal.
Thus, the Fefferman metric $\wh{G}^{\mathrm{F}}$ on $\wh{\mathcal{P}}$ satisfies
\[
\d^* \wh{G}^{\mathrm{F}}={\rm e}^\Upsilon G^{\mathrm{F}}.
\]

\item[(iii)] The conformal metric $G^{\mathrm{F}}$ on $\mathcal{P}$ is also considered by Holland~\cite{H}, and it is proved in {\cite[Section~3.5.4]{H}} that null geodesics of this metric project to conformal geodesics on $\Sigma$. This gives another proof to Theorem~\ref{null-geod-proj}.
\end{itemize}
\end{Remark}

In Section~\ref{lift}, we defined a canonical lift of a regular curve $x(t)$ on $\Sigma$ to a curve $\tilde x(t)=(x(t), [\zeta(t)])$ on $M$ which becomes a chain if and only if $x(t)$ is a conformal geodesic. This lift corresponds to the curve $(x(t), \dot x(t)/|\dot x(t)|)$ on $S\Sigma$ via the isomorphism~\eqref{M-isom}. Thus, if we regard $S\Sigma$ as the twistor CR manifold,
Theorem~\ref{canonical-lift} yields the following.

\begin{Theorem}\label{chain-conf-geod2}
For a regular curve $x(t)$ on $(\Sigma, g)$, the curve $\tilde x(t)=(x(t), \dot x(t)/|\dot x(t)|)$ on $S\Sigma$ is a~chain if and only if $x(t)$ is a conformal geodesic of $[g]$. Moreover, for a conformal geodesic $x(t)$, $a\le t\le b$, the curve $\tilde x(t)$ is a unique chain through $(x(a), \dot x(a)/|\dot x(a)|)$ which projects to~$x(t)$.
\end{Theorem}

We can compete the proof of Theorem~\ref{chain-projection}\,(2) by making use of the right $\mathrm{SO}(3)$-invariance of~$G^{\mathrm{F}}$.

\begin{proof}[Proof of Theorem~\ref{chain-projection}\,(2)]
Let $x(t)$, $a\le t\le b$, be a conformal geodesic on $\Sigma$ and $\eta_0\in S_{x(a)}\Sigma$. Take a null geodesic $\gamma(t)$ on $\bigl(\mathcal{P}, G^{\mathrm{F}}\bigr)$ which projects to the chain $\tilde x(t)=(x(t), \dot x(t)/|\dot x(t)|)$.
We choose $A\in \mathrm{SO}(3)$ so that the third basis vector in $R_A(\gamma(a))$ coincides with $\eta_0$, and set $\gamma'(t):=R_A(\gamma(t))$. Since $G^{\mathrm{F}}$ is right-invariant, $\gamma'$ is also a null geodesic. Projecting $\gamma'$ to $S\Sigma$, we obtain a chain ($g(\eta_0, \dot x(a))\neq0$) or a null chain
 ($g(\eta_0, \dot x(a))=0$) through $\eta_0$. The uniqueness also follows from that of $\tilde x(t)$ and the fact that a lift of chain to a null geodesic is uniquely determined by the initial value.
\end{proof}
